\begin{pastexam}[label={px:df3-111-ds-03}]{地方三等資訊處理 111 年 第三題}
（申論；25 分）

回顧二元樹結構，其為 $m$ 路樹（$m$-ary Trees，亦稱多元樹、$m$ 元樹）的一個特例，請回答下列相關問題：
\begin{itemize}[leftmargin=2.2em, nosep]
\item[㈠] 給出 $m$ 路樹的定義。（5 分）
\item[㈡] 若用陣列來表示一個 $m$ 路樹，請說明如何利用陣列的索引值來表示節點間的親子連結關係（意即，假設陣列索引
起始值為 $0$，若節點 $v$ 在陣列的第 $i$ 個位置，節點 $v$ 的第 $c$ 個子節點的位置為何？另一方面，節點 $v$ 的
parent 位置為何？）？（10 分）
\item[㈢] 基於此 $m$ 路樹結構及二元搜尋樹（Binary Search Tree）的概念，我們可以定義出一個多元搜尋樹。當 $m$=4 的
時候，可以稱此搜尋樹為四元搜尋樹。請給出 (2,4)-樹 ((2,4)-tree) 的定義並比較與四元搜尋樹的差異。（10 分）
\end{itemize}
\end{pastexam}

\begin{pastsoln}
答案：㈡ 第 $c$ 個子節點（$c = 1, \dots, m$）在 $m i + c$，父節點在 $\lfloor (i - 1)/m \rfloor$（$i \ge 1$）。
㈢ (2,4)-樹是每個內部節點有 $2$ 到 $4$ 個子節點、所有葉子同一層的四元搜尋樹；四元搜尋樹沒有這兩條限制，可能
長歪。

（㈠㈡ 屬資結書第~7~章，㈢ 屬第~11~章，以㈢為主。）

\textbf{㈠ 定義。} $m$ 路樹（$m$-ary tree）是有限個節點的集合，它是空的，或者由一個根節點加上 $m$ 棵\textbf{有次序}的
子樹 $T_1, T_2, \dots, T_m$ 組成，每一棵子樹也是 $m$ 路樹（可以是空的）。 換句話說，每個節點最多有 $m$ 個子節點，
而且子節點的位置（第幾個）有意義。 $m = 2$ 就是二元樹：左子樹是 $T_1$、右子樹是 $T_2$（資結書第~7.1~節）。

\textbf{㈡ 陣列表示。} 和二元樹的陣列表示一樣（資結書第~7.2~節），把樹想成補滿的完整 $m$ 路樹，由上而下、由左而右
從 $0$ 開始編號，節點放在自己編號的位置，空的位置留空。
\begin{itemize}
\item \textbf{子節點}：編號比 $i$ 小的節點有 $i$ 個，它們的子節點都排在 $i$ 的子節點前面，共 $m i$ 個位置；再加上
根佔的位置 $0$，所以 $v$ 的第 $c$ 個子節點在
\[
m i + c,\qquad c = 1, 2, \dots, m.
\]
\item \textbf{父節點}：反過來，位置 $j \ge 1$ 可以寫成 $j = m i + c$，$1 \le c \le m$，也就是 $j - 1 = m i + (c - 1)$，
$0 \le c - 1 \le m - 1$，所以
\[
\mathit{parent}(j) = \left\lfloor \frac{j - 1}{m} \right\rfloor,\qquad c = j - m\,\mathit{parent}(j).
\]
根（位置 $0$）沒有父節點。
\end{itemize}
$m = 2$ 時就是熟悉的 $2i + 1$、$2i + 2$ 和 $\lfloor (i - 1)/2 \rfloor$。 （$m = 3$ 的例子見
歷屆試題~\ref{px:gk3-100-ds-06}。）

\textbf{㈢ (2,4)-樹的定義。} 先說四元搜尋樹（4-way search tree，資結書第~11.1~節的 $m$ 路搜尋樹取 $m = 4$）：每個
節點最多 $4$ 棵子樹；有 $d$ 棵子樹的節點存 $d - 1$ 個由小到大的鍵 $a_1 < \cdots < a_{d-1}$，第 $j$ 棵子樹的鍵都在
$a_{j-1}$ 和 $a_j$ 之間（兩端視為 $-\infty$、$+\infty$）。 (2,4)-樹（也叫 2-3-4 樹，資結書第~11.1.5~節）是滿足下面
兩條額外規定的四元搜尋樹：
\begin{enumerate}
\item \textbf{大小性質}：每個內部節點有 $2$、$3$ 或 $4$ 個子節點，也就是存 $1$ 到 $3$ 個鍵（根也至少 $1$ 個鍵）。
\item \textbf{深度性質}：所有葉子（或者說所有外部的空子樹）都在同一層。
\end{enumerate}

\textbf{差異。}
\begin{center}
\small
\begin{tabular}{p{0.16\linewidth}p{0.36\linewidth}p{0.36\linewidth}}
\toprule
 & 四元搜尋樹 & (2,4)-樹 \\
\midrule
每個節點 & 最多 $4$ 個子節點，下限沒規定，可以只有一個鍵、一個非空子樹 & $2$ 到 $4$ 個子節點，不能太空 \\
葉子的深度 & 沒規定 & 全部相同 \\
高度 & 最壞 $n$ 個鍵排成一條鏈，高度 $\Theta(n)$ & $\log_4(n+1) \le h \le \lg(n+1)$，$\Theta(\log n)$ \\
搜尋、插入、刪除 & 最壞 $O(n)$ & $O(\log n)$ \\
維持方法 & 插入只是放進去，不調整 & 插入時節點滿了（$4$ 個鍵）就分裂，刪除時鍵太少就向兄弟借或合併 \\
\bottomrule
\end{tabular}
\end{center}
所以 (2,4)-樹是四元搜尋樹的一種，多了「夠滿」和「葉子同層」兩條規定，換來高度一定是 $O(\log n)$。 它就是 $m = 4$
的 B 樹（資結書第~11.2~節），也和紅黑樹一一對應（資結書第~10.3~節）。
\end{pastsoln}
